\documentclass[11pt]{article}
\usepackage{amsmath}
\usepackage[hidelinks]{hyperref}
\title{Two corridor density demands with a notice received in motion}
\date{27 November 2026}
\begin{document}
\maketitle
\paragraph{Audit.}
This exact continuation of the corridor family removes the
mandatory rest at the notice cut in PR~710, for two contiguous
blocks only. All robots start and end at rest and have the
same acceleration bound \(a>0\) and speed bound \(v>0\).
Every adjacent gap remains in \([g,R]\). Each block contains
\(m\) gaps initially of size \(g\). The two worlds require an
extra terminal span \(\delta>0\) in the left or the right block.
There is adequate extent at both ends. All agents receive an
identical identifying notice at time \(\tau\in[0,T]\), with
no distinguishing evidence before it. This is a condition on
actual common receipt, not a way to obtain a timely notice.
Use the rest-to-rest displacement budget \(B(H)\) of PR~709
for a phase of length \(H\). Physical feasibility requires
\(\delta\le m(R-g)\).

\paragraph{Exact moving-cut deadline.}
Subject to the geometric cap, the necessary and sufficient
condition is
\[
                       \delta\le B(T)+B(T-\tau).
\]
For necessity, let \(L_i,M_i,C_i\) be the final displacements
of the left, middle and right boundaries in world \(i\).
The left demand requires \(M_{\rm L}-L_{\rm L}\ge\delta\),
and the right requires \(C_{\rm R}-M_{\rm R}\ge\delta\).
Adding them gives the bound: the two outside terms together
contribute at most \(2B(T)\). The half-difference of the
two middle trajectories has zero position and velocity at
\(\tau\), since their local motion was common through that
cut. It ends at rest and still satisfies both actuator bounds
by convexity. Consequently \((M_{\rm L}-M_{\rm R})/2\le
B(T-\tau)\). This bounds all admissible causal trajectory
pairs, not only those that choose a symmetric prefix.

For attainment let \(F(t)\) be integrated velocity
\(\min(at,a(T-t),v)\) from 0 to \(t\).
Let \(G(t)\) be the same integrated envelope for the shorter
horizon \(T-\tau\), starting from zero at \(\tau\) and zero
before then. Its speed at any common instant is no greater
than the speed of \(F\), since the remaining time agrees and
its elapsed time is smaller. Hence \(0\le G(t)\le F(t)\).
Put \(\alpha=\delta/(B(T)+B(T-\tau))\), which is at most 1 in
the feasible case. The outside boundaries move by
\(-\alpha F(t)\) and \(+\alpha F(t)\) in both worlds.
The middle boundary moves by \(+\alpha G(t)\) for the left
demand and \(-\alpha G(t)\) for the right demand. The actual
motions are known before the notice except for that sign;
the middle stays still until the receipt. The two span
increments are \(\alpha(F+G)\) and \(\alpha(F-G)\), both
in \([0,\delta]\). Linearly interpolate displacements within
each block. Convexity gives each interior robot the same
actuator limits, and all hard clearances and geometric range
bounds hold at every instant. Each robot comes to rest at
\(T\). For \(T=0\) no positive demand is feasible.

\paragraph{Matched controls.}
A fixed blind motion must handle both demands; applying the
two static rows gives exact capacity \(B(T)\). Waiting to
move until after the notice gives \(2B(T-\tau)\). A mandatory
rest cut instead gives \(B(\tau)+2B(T-\tau)\): each block
can add at most \(2B(T-\tau)\) after the cut, and whatever
shortfall remains is required in both blocks at the common
cut. The pre-cut ends can provide at most \(2B(\tau)\) of
total extra span; shared preparation split equally attains
this limit when the cap permits. At \(a=2,T=6/5,\tau=2/5\)
with no binding speed cap, the moving capacity is \(26/25\),
forced rest and blind are \(18/25\), and waiting is \(16/25\).
Only the moving informed policy can meet the unit demand in
that comparison. With \(T=2,\delta=9/4,a=2\), the latest
informative notice is \(\tau=2-\sqrt{1/2}\) under a nonbinding
speed cap and adequate geometric span. A same-information
direct scenario-tree controller cannot improve the half-difference
bound and can use the same witness.

\paragraph{Verification and scope.}
The closed piecewise primitives agree with independent
trapezoidal accumulation at 133,600 exact dyadic sample cuts
across 3,040 parameter combinations; the analytic proof is what
covers intermediate times. Four Lean lemmas cover the real
algebra of the bounds. The extra rational comparison is
verified in integer-scaled time and displacement units.
No field movement, communication latency, private partial
notice, more than two demanded alternatives, or R6 closure
is established. In particular this does not solve general
moving-cut code design.
\end{document}
